% GENERATED FILE. Edit canonical lessons, then run npm run generate:books.
% canonical-source-hash: fnv1a64:cae6625ef95abb4b
% canonical-lessons: BN-C306-otithi, BN-C306-nimontron-kora, BN-C306-bhoj, BN-C306-upohar, BN-C306-apyayon-kora

\chapter{Guests and a Feast}
\label{ch:bn-guests-and-a-feast}
\hlchaptermodality{306}

\begin{chapteropening}
\textbf{By the end of this chapter:} \emph{I can talk about invitations, guests, gifts and a feast.}

Complete the last lesson of chapter 306: I can talk about invitations, guests, gifts and a feast.
\end{chapteropening}

\section[\texorpdfstring{\bn{অতিথি} (ôtithi)}{অতিথি (ôtithi)}]{\bn{অতিথি} (ôtithi) — a guest}

\label{lesson:BN-C306-otithi}



\begin{quote}
Before the new one: say the Bengali for to glitter, then the Bengali for a sky lantern.
\end{quote}

\subsection*{You'll want to know: \bn{অতিথি}}
\textbf{\bn{অতিথি}} — \emph{ôtithi} — \textquotedblleft{}a guest\textquotedblright{}.

\begin{quote}
\bn{পুজোর} \bn{দিন} \bn{বাড়িতে} \bn{অনেক} \bn{অতিথি} \bn{এলেন।}
\end{quote}

\begin{quote}
\emph{pujor din bāṛite ônek ôtithi elen.}
\end{quote}

\begin{quote}
Many guests came to the house on the day of the puja.
\end{quote}

\begin{cousinweb}
\textbf{\bn{অতিথি}} (\emph{ôtithi}) is Sanskrit \emph{atithi}, 'a guest, one owed hospitality' (Monier-Williams), which records an old explanation: \emph{a-tithi}, one with no fixed day for coming.
\end{cousinweb}

\begin{grammarlens}[title={What you've built}]
One more word for telling what happened.
\end{grammarlens}

\subsection*{Your turn}
\begin{itemize}
\raggedright
  \item \emph{Say it:} \emph{ôtithi}
  \item \emph{Say it:} \emph{ôtithi}, once more
  \item \emph{Say it:} \emph{pujor din bāṛite ônek ôtithi elen.} — Many guests came to the house on the day of the puja.
  \item \emph{Recall:} say the Bengali for to glitter, then the Bengali for a sky lantern, then say \emph{ôtithi} again
\end{itemize}

\begin{tcolorbox}[breakable,colback=teal!4,colframe=teal!35!black,title={Before you move on}]
What does \bn{অতিথি} mean? (\textquotedblleft{}A guest\textquotedblright{}.) Say it once more.
\end{tcolorbox}
\section[\texorpdfstring{\bn{নিমন্ত্রণ} \bn{করা} (nimôntrôn kôrā)}{নিমন্ত্রণ করা (nimôntrôn kôrā)}]{\bn{নিমন্ত্রণ} \bn{করা} (nimôntrôn kôrā) — to invite}

\label{lesson:BN-C306-nimontron-kora}



\begin{quote}
Before the new one: say the Bengali for a sky lantern, then the Bengali for a guest.
\end{quote}

\subsection*{You'll want to know: \bn{নিমন্ত্রণ} \bn{করা}}
\textbf{\bn{নিমন্ত্রণ} \bn{করা}} — \emph{nimôntrôn kôrā} — \textquotedblleft{}to invite\textquotedblright{}.

\begin{quote}
\bn{মামা} \bn{আমাদের} \bn{তাঁর} \bn{মেয়ের} \bn{বিয়েতে} \bn{নিমন্ত্রণ} \bn{করলেন।}
\end{quote}

\begin{quote}
\emph{māmā āmāder tā̃r meyer biyete nimôntrôn kôrlen.}
\end{quote}

\begin{quote}
My uncle invited us to his daughter's wedding.
\end{quote}

Past (he/she …): \bn{তিনি} \bn{আমাদের} \bn{নিমন্ত্রণ} \bn{করলেন}, \emph{tini āmāder nimôntrôn kôrlen}.

\begin{cousinweb}
\textbf{\bn{নিমন্ত্রণ} \bn{করা}} (\emph{nimôntrôn kôrā}) puts Sanskrit \emph{nimantraṇa}, 'an invitation, a summons' (Monier-Williams), before \bn{করা}.
\end{cousinweb}

\begin{grammarlens}[title={What you've built}]
One more word for telling what happened.
\end{grammarlens}

\subsection*{Your turn}
\begin{itemize}
\raggedright
  \item \emph{Say it:} \emph{nimôntrôn kôrā}
  \item \emph{Say it:} \emph{nimôntrôn kôrā}, once more
  \item \emph{Say it:} \emph{māmā āmāder tā̃r meyer biyete nimôntrôn kôrlen.} — My uncle invited us to his daughter's wedding.
  \item \emph{Recall:} say the Bengali for a sky lantern, then the Bengali for a guest, then say \emph{nimôntrôn kôrā} again
\end{itemize}

\begin{tcolorbox}[breakable,colback=teal!4,colframe=teal!35!black,title={Before you move on}]
What does \bn{নিমন্ত্রণ} \bn{করা} mean? (\textquotedblleft{}To invite\textquotedblright{}.) Say it once more.
\end{tcolorbox}
\section[\texorpdfstring{\bn{ভোজ} (bhoj)}{ভোজ (bhoj)}]{\bn{ভোজ} (bhoj) — a feast}

\label{lesson:BN-C306-bhoj}



\begin{quote}
Before the new one: say the Bengali for a guest, then the Bengali for to invite.
\end{quote}

\subsection*{You'll want to know: \bn{ভোজ}}
\textbf{\bn{ভোজ}} — \emph{bhoj} — \textquotedblleft{}a feast\textquotedblright{}.

\begin{quote}
\bn{বিয়ের} \bn{রাতে} \bn{একটা} \bn{বড়} \bn{ভোজ} \bn{হল।}
\end{quote}

\begin{quote}
\emph{biyer rāte ekṭā bôṛo bhoj hôlô.}
\end{quote}

\begin{quote}
There was a big feast on the night of the wedding.
\end{quote}

\begin{cousinweb}
\textbf{\bn{ভোজ}} (\emph{bhoj}) comes down from Sanskrit \emph{bhojya}, 'to be eaten; a meal' (CDIAL 9631, which lists Bengali \emph{bhoj}, 'a feast').
\end{cousinweb}

\begin{grammarlens}[title={What you've built}]
One more word for telling what happened.
\end{grammarlens}

\subsection*{Your turn}
\begin{itemize}
\raggedright
  \item \emph{Say it:} \emph{bhoj}
  \item \emph{Say it:} \emph{bhoj}, once more
  \item \emph{Say it:} \emph{biyer rāte ekṭā bôṛo bhoj hôlô.} — There was a big feast on the night of the wedding.
  \item \emph{Recall:} say the Bengali for a guest, then the Bengali for to invite, then say \emph{bhoj} again
\end{itemize}

\begin{tcolorbox}[breakable,colback=teal!4,colframe=teal!35!black,title={Before you move on}]
What does \bn{ভোজ} mean? (\textquotedblleft{}A feast\textquotedblright{}.) Say it once more.
\end{tcolorbox}
\section[\texorpdfstring{\bn{উপহার} (upôhār)}{উপহার (upôhār)}]{\bn{উপহার} (upôhār) — a gift; a present}

\label{lesson:BN-C306-upohar}



\begin{quote}
Before the new one: say the Bengali for to invite, then the Bengali for a feast.
\end{quote}

\subsection*{You'll want to know: \bn{উপহার}}
\textbf{\bn{উপহার}} — \emph{upôhār} — \textquotedblleft{}a gift; a present\textquotedblright{}.

\begin{quote}
\bn{জন্মদিনে} \bn{বন্ধুরা} \bn{আমাকে} \bn{একটা} \bn{বই} \bn{উপহার} \bn{দিল।}
\end{quote}

\begin{quote}
\emph{jônmôdine bôndhurā āmāke ekṭā bôi upôhār dilô.}
\end{quote}

\begin{quote}
My friends gave me a book as a birthday present.
\end{quote}

\begin{cousinweb}
\textbf{\bn{উপহার}} (\emph{upôhār}) is Sanskrit \emph{upahāra}, an offering to a god or a present to a superior (Monier-Williams); it has nothing to do with \textbf{\bn{হার}}, a necklace. \bn{উপহার} \bn{দেওয়া}, to give a present.
\end{cousinweb}

\begin{grammarlens}[title={What you've built}]
One more word for telling what happened.
\end{grammarlens}

\subsection*{Your turn}
\begin{itemize}
\raggedright
  \item \emph{Say it:} \emph{upôhār}
  \item \emph{Say it:} \emph{upôhār}, once more
  \item \emph{Say it:} \emph{jônmôdine bôndhurā āmāke ekṭā bôi upôhār dilô.} — My friends gave me a book as a birthday present.
  \item \emph{Recall:} say the Bengali for to invite, then the Bengali for a feast, then say \emph{upôhār} again
\end{itemize}

\begin{tcolorbox}[breakable,colback=teal!4,colframe=teal!35!black,title={Before you move on}]
What does \bn{উপহার} mean? (\textquotedblleft{}A gift; a present\textquotedblright{}.) Say it once more.
\end{tcolorbox}
\section[\texorpdfstring{\bn{আপ্যায়ন} \bn{করা} (āpyāyôn kôrā)}{আপ্যায়ন করা (āpyāyôn kôrā)}]{\bn{আপ্যায়ন} \bn{করা} (āpyāyôn kôrā) — to entertain guests; to treat hospitably}

\label{lesson:BN-C306-apyayon-kora}



\begin{quote}
Before the new one: say the Bengali for a feast, then the Bengali for a gift.
\end{quote}

\subsection*{You'll want to know: \bn{আপ্যায়ন} \bn{করা}}
\textbf{\bn{আপ্যায়ন} \bn{করা}} — \emph{āpyāyôn kôrā} — \textquotedblleft{}to entertain guests; to treat hospitably\textquotedblright{}.

\begin{quote}
\bn{বাবা} \bn{অতিথিদের} \bn{চা} \bn{আর} \bn{সন্দেশ} \bn{দিয়ে} \bn{আপ্যায়ন} \bn{করলেন।}
\end{quote}

\begin{quote}
\emph{bābā ôtithider chā ār sôndesh diye āpyāyôn kôrlen.}
\end{quote}

\begin{quote}
Father welcomed the guests with tea and sweets.
\end{quote}

Past (he/she …): \bn{তিনি} \bn{অতিথিদের} \bn{আপ্যায়ন} \bn{করলেন}, \emph{tini ôtithider āpyāyôn kôrlen}.

\begin{cousinweb}
\textbf{\bn{আপ্যায়ন} \bn{করা}} (\emph{āpyāyôn kôrā}) puts Sanskrit \emph{āpyāyana}, 'causing fullness, gladdening' (Monier-Williams), before \bn{করা}: to fill your guests up and make them glad.
\end{cousinweb}

\begin{grammarlens}[title={What you've built}]
One more word for telling what happened.
\end{grammarlens}

\subsection*{Your turn}
\begin{itemize}
\raggedright
  \item \emph{Say it:} \emph{āpyāyôn kôrā}
  \item \emph{Say it:} \emph{āpyāyôn kôrā}, once more
  \item \emph{Say it:} \emph{bābā ôtithider chā ār sôndesh diye āpyāyôn kôrlen.} — Father welcomed the guests with tea and sweets.
  \item \emph{Recall:} say the Bengali for a feast, then the Bengali for a gift, then say \emph{āpyāyôn kôrā} again
\end{itemize}

\begin{tcolorbox}[breakable,colback=teal!4,colframe=teal!35!black,title={Before you move on}]
What does \bn{আপ্যায়ন} \bn{করা} mean? (\textquotedblleft{}To entertain guests; to treat hospitably\textquotedblright{}.) Say it once more.
\end{tcolorbox}
